\chapter{Learning compact genomic representations}
\label{chap:representations}
\label{chap:latent_representations}

After defining how genomic data are prepared and separated, the next
question is what kind of representation should be learned from them.
The original post-QC genotype vector still contains more than forty
thousand markers, while the population structure of interest may
occupy a much smaller effective space.

This chapter develops two different answers to that question. The
first explicitly combines genotype reconstruction and breed prediction
in a variational architecture. The second removes breed supervision
from representation learning and asks whether a self-supervised
contrastive objective can organize individuals in a very small
hyperspherical latent space.

The purpose of the chapter is to explain the models and the reasoning
behind their design. Their quantitative comparison is deliberately
deferred to Chapter~\ref{chap:evaluation}.

\section{Why learn a latent genomic space?}
\label{sec:why_latent_genomics}

Using the complete genomic vector directly has an intuitive advantage:
no marker is discarded before learning. At the same time, a
representation with tens of thousands of coordinates is difficult to
inspect and may contain substantial redundancy due to LD.

A latent space changes the problem. Instead of representing an
individual by every measured locus, the model is asked to summarize
the input through a much smaller vector. The usefulness of this
operation depends on what is preserved by the compression.

A reconstruction-based model is encouraged to retain information that
describes the original genotype profile. A supervised classifier is
encouraged to retain information that distinguishes breeds. A
contrastive model, in turn, is encouraged to preserve information that
remains stable across different views of the same individual while
separating different individuals.

These objectives are not equivalent. Comparing them provides a way to
investigate whether direct breed supervision is necessary to obtain a
discriminative genomic representation, or whether breed structure also
emerges from a self-supervised organization of genotype profiles. Any
learning procedure encodes assumptions about the structure of the data,
i.e. an inductive bias \cite{mitchell1980need}, and the two objectives
below differ precisely in the assumptions they impose on a useful
genomic representation.

\section{A variational approach to genotype representation}
\label{sec:vmgp_architecture}

The first model follows the variational autoencoder (VAE) paradigm
\cite{kingma2013auto, rezende2014stochastic} and is inspired by VMGP
\cite{zhao2025vmgp}. The original VMGP formulation was developed for
plant genomic prediction. Here, its architectural ideas are adapted to
livestock breed classification.

Let $P$ denote the number of markers retained by the preprocessor for
the current training partition. The variational encoder receives a
dosage vector $\mathbf{x}\in\mathbb{R}^{P}$ and processes it through

\begin{equation}
    P
    \rightarrow
    2048
    \rightarrow
    512.
\end{equation}

Each hidden stage uses a linear transformation followed by Batch
Normalization, LeakyReLU with negative slope $0.2$, and dropout with
probability $0.1$. Two projections then produce the parameters of a
diagonal Gaussian posterior,

\begin{equation}
    \boldsymbol{\mu}(\mathbf{x}),
    \log\boldsymbol{\sigma}^{2}(\mathbf{x})
    \in\mathbb{R}^{d}.
\end{equation}

\subsection{Sampling and reconstruction}

During training, a latent vector is sampled using the
reparameterization trick,

\begin{equation}
    \mathbf{z}
    =
    \boldsymbol{\mu}
    +
    \exp
    \left(
        \frac{1}{2}\log\boldsymbol{\sigma}^{2}
    \right)
    \odot
    \boldsymbol{\epsilon},
    \qquad
    \boldsymbol{\epsilon}
    \sim
    \mathcal{N}(\mathbf{0},\mathbf{I}).
    \label{eq:reparameterization}
\end{equation}

The decoder mirrors the encoder,

\begin{equation}
    d
    \rightarrow
    512
    \rightarrow
    2048
    \rightarrow
    P,
\end{equation}

and reconstructs the original allele-dosage vector.

At evaluation time, the stochastic sampling step is removed and the
deterministic posterior mean $\boldsymbol{\mu}$ is used as the
representation.

\subsection{Breed prediction from the latent mean}

Reconstruction alone does not guarantee that the latent space will
organize individuals in a way that is convenient for breed
classification. A supervised head therefore receives the
posterior mean and maps it through

\begin{equation}
    d
    \rightarrow
    128
    \rightarrow
    C,
\end{equation}

where $C$ is the number of breeds. The hidden layer uses Batch
Normalization, LeakyReLU with slope $0.2$, and dropout with
probability $0.3$. The final linear layer produces breed logits.

Using the posterior mean rather than a stochastic sample for breed
prediction also gives the classifier a deterministic representation
on which to operate.

\subsection{Joint reconstruction and classification objective}

The reconstruction term combines element-wise mean squared error with
a small KL regularizer,

\begin{equation}
    \mathcal{L}_{\mathrm{rec}}
    =
    \mathcal{L}_{\mathrm{MSE}}
    +
    10^{-4}
    \mathcal{L}_{\mathrm{KL}},
\end{equation}

where

\begin{equation}
    \mathcal{L}_{\mathrm{KL}}
    =
    -\frac{1}{2B}
    \sum_{n=1}^{B}
    \sum_{j=1}^{d}
    \left(
        1
        +
        \log\sigma_{n,j}^{2}
        -
        \mu_{n,j}^{2}
        -
        \sigma_{n,j}^{2}
    \right).
\end{equation}

Breed prediction uses multiclass cross-entropy,

\begin{equation}
    \mathcal{L}_{\mathrm{breed}}
    =
    \operatorname{CE}
    \left(
        \mathbf{l}_{\mathrm{breed}},
        \mathbf{y}
    \right).
\end{equation}

The executable objective is

\begin{equation}
    \mathcal{L}_{\mathrm{VMGP}}
    =
    \alpha\mathcal{L}_{\mathrm{rec}}
    +
    (1-\alpha)\mathcal{L}_{\mathrm{breed}},
    \qquad
    \alpha=0.5.
    \label{eq:vmgp_total_loss}
\end{equation}

The KL coefficient is therefore fixed at $10^{-4}$ inside the
reconstruction term. The implementation is not treated as an active
$\beta$-VAE, and no auxiliary continent, casein, or aptitude
prediction head contributes to the training loss.

\subsection{How large should the variational representation be?}

A latent representation should be small enough to provide meaningful
compression but large enough to preserve the structure required by the
task. Rather than selecting a dimensionality by intuition, four
candidate sizes are considered during development,

\begin{equation}
    d\in\{32,64,96,128\}.
\end{equation}

The choice among them is made only from the three cross-validation folds
using the selection rule defined in Section~\ref{sec:rq1_protocol}.
The purpose of separating the candidate definition from its evaluation
is to avoid rewriting the architecture around the outcome of the
final test.

Optimization uses Adam with a learning rate $10^{-4}$ and a batch size
of $64$. Training uses mixed precision. The development ceiling is
$200$ epochs, while EarlyStopping monitors validation loss with
patience $15$. The checkpoint with minimum validation loss is restored
before fold-level classification metrics are computed.

\begin{figure}[htbp]
\centering
\begin{tikzpicture}[
    >=Latex,
    font=\small,
    box/.style={draw, semithick, rounded corners=2pt, align=center,
                minimum height=10mm, minimum width=42mm, inner sep=4pt, fill=black!2},
    key/.style={box, fill=black!7},
    small/.style={draw, semithick, rounded corners=2pt, align=center,
                  minimum height=9mm, minimum width=29mm, inner sep=3pt, fill=black!2},
    arr/.style={->, semithick}
]
\node[key] (x) at (0,4.6) {Encoded SNP profile $\mathbf{x}$};
\node[box] (enc) at (0,3.2) {Encoder $q_\phi$};
\node[small] (mu) at (-2.2,1.75) {$\boldsymbol{\mu}(\mathbf{x})$};
\node[small] (lv) at (2.2,1.75) {$\log\boldsymbol{\sigma}^{2}(\mathbf{x})$};
\node[key] (z) at (0,0.15) {Reparameterization $\rightarrow$ latent $\mathbf{z}$};
\node[box] (dec) at (-3.0,-1.35) {Decoder};
\node[box] (head) at (3.0,-1.35) {Breed prediction head};
\node[small] (rec) at (-3.0,-2.8) {Reconstruction};
\node[small] (prob) at (3.0,-2.8) {Breed probabilities};
\node[key] (loss) at (0,-4.35) {Joint objective $\alpha\mathcal{L}_{\mathrm{rec}}+(1-\alpha)\mathcal{L}_{\mathrm{breed}}$, with $\mathcal{L}_{\mathrm{rec}}=\mathcal{L}_{\mathrm{MSE}}+10^{-4}\mathcal{L}_{\mathrm{KL}}$};

\draw[arr] (x)--(enc); \draw[arr] (enc)--(mu); \draw[arr] (enc)--(lv);
\draw[arr] (mu)--(z); \draw[arr] (lv)--(z);
\draw[arr] (z)--(dec); \draw[arr] (z)--(head);
\draw[arr] (dec)--(rec); \draw[arr] (head)--(prob);
\draw[arr] (rec)--(loss); \draw[arr] (prob)--(loss);
\end{tikzpicture}

\caption{Variational multi-task architecture used as the basis of the breed-aware genomic representation model. The encoder parameterizes a latent posterior regularized by a KL-divergence term; the resulting latent variable supports both genotype reconstruction and breed prediction. Author's adaptation of the VAE formulation \cite{kingma2013auto} and the VMGP multi-task paradigm \cite{zhao2025vmgp}.}
\label{fig:vmgp_architecture}
\end{figure}


\section{A self-supervised contrastive approach}
\label{sec:contrastive_model}

The second model starts from a different question. Instead of telling
the representation learner which breed an individual belongs to, can
a useful genomic space be learned from the relationship between
different views of the same individual?

The implementation follows the paper-faithful contrastive design of
Thor and Nettelblad \cite{thor2025contrastive}. Its core components are
a genomic augmentation operator, a convolutional encoder, normalization
onto a hypersphere, and a centroid N-pair objective. Breed labels do not
enter the contrastive training loss.

\subsection{Genomic views and data augmentation}

The shared preprocessing pipeline produces dosage states $0$, $1$,
and $2$. For contrastive learning, these states are represented
categorically using four channels. The first three correspond to the
biological genotype states, while the fourth is reserved for an
augmentation-only mask state $-1$,

\begin{equation}
    \tilde{g}_{i,j}
    \in
    \{-1,0,1,2\}
    \longrightarrow
    \mathbf{x}_{i,j}
    \in
    \{0,1\}^{4}.
    \label{eq:contrastive_one_hot}
\end{equation}

During training, two independently augmented views of the same
individual are generated,

\begin{equation}
    \tilde{\mathbf{x}}_i^{a},
    \tilde{\mathbf{x}}_i^{b}
    \sim
    \mathcal{T}(\mathbf{x}_i).
    \label{eq:contrastive_views}
\end{equation}

The implemented augmentation combines genotype-state perturbation and
masking. For each generated view, a perturbation probability is
sampled uniformly in the configured interval from $0.01$ to $0.99$.
Homozygous states can be moved toward the heterozygous state
($0\rightarrow1$ and $2\rightarrow1$), while a heterozygous state can
be perturbed toward either homozygous state. A masking probability is
then sampled in the same interval and the selected positions are
assigned to the special state $-1$.

The augmentation is intentionally disabled during deterministic
evaluation. At that point, the fourth channel is zero and the embedding
is derived from the unperturbed post-QC genotype profile.

\subsection{Encoder and hyperspherical embedding}

The categorical tensor has shape $(B,4,P)$ and is processed by two
one-dimensional convolutions,

\begin{equation}
\operatorname{Conv1d}(4,5,3)
\rightarrow
\operatorname{SiLU}
\rightarrow
\operatorname{Conv1d}(5,5,3)
\rightarrow
\operatorname{SiLU},
\end{equation}

using padding one in both layers.

After flattening, the representation passes through

\begin{equation}
    5P
    \rightarrow
    256
    \rightarrow
    256
    \rightarrow
    256
    \rightarrow
    3.
\end{equation}

Each $256$-unit hidden layer uses Batch Normalization and SiLU.

The final vector is normalized,

\begin{equation}
    \mathbf{z}
    =
    \frac{\mathbf{h}}{\|\mathbf{h}\|_2},
    \qquad
    \|\mathbf{z}\|_2=1.
    \label{eq:hyperspherical_normalization}
\end{equation}

The resulting three-dimensional embeddings, therefore, lie on
$\mathbb{S}^{2}\subset\mathbb{R}^{3}$. This extremely small
representation is a deliberate part of the contrastive design rather
than a value selected from the same grid used by the variational
model.

\subsection{Centroid N-pair learning}

For an anchor $\mathbf{z}_i$, a positive view $\mathbf{z}_i^{+}$, and
a negative embedding $\mathbf{z}_{ij}^{-}$, the implemented objective
first defines a centroid

\begin{equation}
    \mathbf{c}_{ij}
    =
    \frac{
        \mathbf{z}_i
        +
        2\mathbf{z}_{ij}^{-}
        +
        \mathbf{z}_i^{+}
    }{4}.
\end{equation}

The three vectors are centered around this centroid and scaled by

\begin{equation}
    r_{ij}
    =
    \sqrt{
        \max
        \left(
            \|\mathbf{z}_i-\mathbf{c}_{ij}\|_2^2,
            \|\mathbf{z}_i^{+}-\mathbf{c}_{ij}\|_2^2,
            \|\mathbf{z}_{ij}^{-}-\mathbf{c}_{ij}\|_2^2
        \right)
        +
        10^{-8}
    }.
\end{equation}

Using the centered and scaled vectors, positive and negative
similarities are

\begin{equation}
    s_{ij}^{+}
    =
    \tilde{\mathbf{z}}_i^\top
    \tilde{\mathbf{z}}_i^{+},
    \qquad
    s_{ij}^{-}
    =
    \tilde{\mathbf{z}}_i^\top
    \tilde{\mathbf{z}}_{ij}^{-}.
\end{equation}

The per-anchor contribution is

\begin{equation}
    \ell_i
    =
    \log
    \left[
        1
        +
        \sum_{j\neq i}
        \max
        \left(
            0,
            \exp(s_{ij}^{-}-s_{ij}^{+})
            -
            e^{-2}
        \right)
    \right].
\end{equation}

The batch loss is the mean of $\ell_i$. This objective is different
from NT-Xent: it does not use a temperature parameter and retains the
centroid-based geometry of the reference implementation.

\subsection{From a self-supervised space to breed prediction}

The contrastive model does not contain a supervised breed head. Once
training for a cross-validation fold is complete, deterministic embeddings
are extracted separately for the training and validation individuals.

A $k$-nearest-neighbor classifier with $k=3$ is then fitted using only
the training embeddings and training breed labels. The validation
embeddings are classified using this fitted $k$-NN.

This two-stage procedure is important for interpretation. A successful
breed prediction does not mean that the contrastive loss was trained
with breed labels. It means that breed structure has become
recoverable from a representation learned through self-supervised
instance relationships.

\subsection{Optimization}

The contrastive model is optimized using Adam with learning rate
$10^{-3}$, $\beta_1=0.9$, and $\beta_2=0.999$. A StepLR scheduler
multiplies the learning rate by $0.99$ every $10$ epochs.

The development ceiling is $5000$ epochs. Early stopping monitors
validation loss with patience $200$ and the checkpoint with minimum
validation loss is restored before the downstream embedding
evaluation.

\begin{figure}[htbp]
\centering
\begin{tikzpicture}[
    >=Latex,
    font=\small,
    box/.style={draw, semithick, rounded corners=2pt, align=center,
                minimum height=10mm, minimum width=37mm, inner sep=4pt, fill=black!2},
    key/.style={box, fill=black!7},
    small/.style={draw, semithick, rounded corners=2pt, align=center,
                  minimum height=9mm, minimum width=31mm, inner sep=3pt, fill=black!2},
    arr/.style={->, semithick}
]
\node[key] (x) at (0,4.65) {Genotype profile $\mathbf{x}_i$};
\node[box] (aa) at (-3.25,3.15) {Augmented view A};
\node[box] (ab) at (3.25,3.15) {Augmented view B};
\node[box] (ea) at (-3.25,1.65) {Shared encoder $f_\theta$};
\node[box] (eb) at (3.25,1.65) {Shared encoder $f_\theta$};
\node[small] (za) at (-3.25,0.15) {$L_2$ norm. $\rightarrow \mathbf{z}_i^{a}$};
\node[small] (zb) at (3.25,0.15) {$L_2$ norm. $\rightarrow \mathbf{z}_i^{b}$};
\node[key] (np) at (0,-1.35) {Centroid N-pair loss (Eqs.~3--6)};
\node[key] (sphere) at (0,-2.90) {Hyperspherical latent space $\mathbb{S}^{2}$};

\draw[arr] (x) -- (aa); \draw[arr] (x) -- (ab);
\draw[arr] (aa) -- (ea); \draw[arr] (ab) -- (eb);
\draw[arr] (ea) -- (za); \draw[arr] (eb) -- (zb);
\draw[<->, semithick] (za) -- node[above]{positive pair} (zb);
\draw[arr] (za) -- (np); \draw[arr] (zb) -- (np);
\draw[arr] (np) -- (sphere);
\end{tikzpicture}

\caption{Self-supervised hyperspherical pipeline. Two independently
augmented views of the same individual are encoded by a shared
convolutional encoder and $L_2$-normalized onto the unit sphere. The
representation is optimized with the centroid N-pair objective, without
a projection head, without NT-Xent temperature and without breed labels
\cite{thor2025contrastive,wang2020understanding}.}
\label{fig:contrastive_architecture}
\end{figure}


\section{Two different views of the same genomic problem}
\label{sec:two_views_genomic_problem}

The two models deliberately solve different learning problems.

The variational architecture receives direct information about breed
identity and balances this supervised objective with reconstruction of
the genotype profile. Its latent variables are therefore shaped by
both population discrimination and input reconstruction.

The contrastive model receives no breed labels while learning its
representation. It is instead constrained by augmentation invariance
and by the geometry induced by positive and negative relationships.

A direct comparison of their classification metrics is meaningful
because both are evaluated on the same development individuals and the
same train-fitted genomic preprocessing. At the same time, the
comparison should not be reduced to the claim that one neural
architecture is universally better. The objectives answer different
representation-learning questions, and the evaluation chapter
considers both predictive performance and the properties of the
resulting latent spaces.

\section{Final development training}
\label{sec:final_training_protocol}

Cross-validation is used for development and configuration selection.
Once a configuration has been fixed, the purpose of final training is
different: all development individuals should contribute to the model
that will be evaluated on the held-out test.

For a selected configuration, let $e_1^{*}$, $e_2^{*}$, and
$e_3^{*}$ denote the one-based epochs of the minimum-validation-loss
checkpoints observed in the three cross-validation folds. The final
training duration is defined as

\begin{equation}
    E_{\mathrm{final}}
    =
    \operatorname{median}
    \left(
        e_1^{*},
        e_2^{*},
        e_3^{*}
    \right).
    \label{eq:final_epoch_budget}
\end{equation}

A new genomic preprocessor is fitted from scratch on the complete
development set. The selected model is then trained using all
development individuals for exactly $E_{\mathrm{final}}$ completed
epochs.

There is no validation split, no early stopping, and no checkpoint
selection based on held-out outcomes during this stage. Only after
final training has finished is the held-out test transformed with the
already fitted development preprocessor and evaluated once.

For the contrastive model, the final $k$-NN is fitted on the
deterministic embeddings of the complete development set and then
applied to the held-out test embeddings.

The numerical values selected during development are intentionally
reported with the experimental evaluation rather than embedded in the
architectural description of this chapter.
